\section{Key Takeaway for the Future of SoTS}\label{sec:literature-review-key-takeaways}
This section synthesizes the findings of the literature review to identify the central gaps that limit current SoTS research and practice. While the results in \secref{sec:literature-review-results} describe the combination of DT and SoS, the analysis here interprets what those results imply for the future development of SoTS. Each takeaway highlights a recurring issue observed across the surveyed studies and explains why addressing it is essential for enabling SoTS to operate in realistic, dynamic environments.


%------------------------------------------------------------
\subsection*{The Emerging Reliability Gap in SoTS}

Taken together, the results of the review point to a consistent pattern: SoTS assume that the data needed for their operation will be available when required, yet provide few mechanisms to continue operating effectively when it is not. The motivations identified in \secref{subsec:literature-review-results-motivations} all assume that systems will have access to the data required to perform their intended functions. Likewise, the architectural structures in \secref{subsec:literature-review-results-architecture} depend on relationships between coordinating and local twins that must remain aligned for the broader system to achieve its objectives. The most prominent SoS dimensions, distribution and independence, highlighted in \secref{subsec:literature-review-results-sos-characteristics}, further emphasize that SoTS operate across dispersed, autonomous constituents whose behaviour cannot be centrally enforced. Together, these findings characterize an operational setting in which maintaining system functionality is inherently challenging, and where existing work provides limited support for ensuring reliable system behaviour.


Reliability is mentioned frequently in the literature, but as shown in \secref{subsec:literature-review-results-reliability}, it is rarely treated as a property that must be preserved under changing or uncertain conditions. Instead, most approaches implicitly assume that systems will remain synchronized, communication will remain available, and information will arrive when needed. This framing largely ignores reliability as a characteristic of ongoing system behaviour that must be actively maintained. In SoTS, however, behaviour evolves continuously and depends on the quality and consistency of incoming information. Under such conditions, correctness may deteriorate gradually as inputs become incomplete, delayed, or inconsistent, even though the system remains operational.


The limited attention to dynamic SoS characteristics such as reconfiguration and evolution, again evident in \secref{subsec:literature-review-results-sos-characteristics}, indicates that many studies conceptualize SoTS as operating in stable settings despite the inherent instability of distributed and autonomous systems. The uneven adoption of standards described in \secref{subsec:literature-review-results-standards} shows progress toward interoperability, but provides little assurance that correct behaviour will be maintained at runtime when expected interactions or observations are disrupted. As a result, existing approaches offer limited means to reason about correctness over time.

Considered as a whole, the literature shows that although the structural and conceptual foundations of SoTS are becoming clearer, the operational capabilities needed to maintain reliable functioning are not yet developed. The primary goals that motivate SoTS, enhanced coordination, shared situational awareness, and informed decision-making, depend on system behaviour that remains correct and continuous even as conditions change. Yet current approaches provide limited support for preserving such behaviour when systems reconfigure, when the data they depend on is disrupted, or when operating environments become unstable. This gap represents a central limitation that must be addressed for SoTS to progress beyond prototype demonstrations and operate reliably in practice.

\begin{recoframe}{Recommendation}
Treat reliability as a primary architectural concern in SoTS, as current approaches depend on stable interactions across distributed constituents without providing mechanisms to ensure them. Future work should introduce explicit support for maintaining reliable behaviour under changing conditions so that SoTS can operate beyond controlled prototype environments.
\end{recoframe}

% Taken together, the results of the review point to a consistent pattern: SoTS assume that the data needed for their operation will be available when required, yet provide few mechanisms to continue operating effectively when it is not. The motivations identified in \secref{subsec:literature-review-results-motivations} all assume that systems will have access to the data required to perform their intended functions. Likewise, the architectural structures in \secref{subsec:literature-review-results-architecture} depend on relationships between coordinating and local twins that must remain aligned for the broader system to achieve its objectives. The most prominent SoS dimensions, distribution and independence, highlighted in \secref{subsec:literature-review-results-sos-characteristics}, further emphasize that SoTS operate across dispersed, autonomous constituents whose behaviour cannot be centrally enforced. Together, these findings characterize an operational setting in which maintaining system functionality is inherently challenging, and where existing work provides limited support for ensuring reliable system behaviour.


% Reliability is mentioned frequently, but as shown in \secref{subsec:literature-review-results-reliability}, it is rarely treated as a property that must be preserved under changing or uncertain conditions. Instead, most approaches implicitly assume that systems will remain synchronized, communication will remain available, and information will arrive when needed. The limited attention to dynamic SoS characteristics such as reconfiguration and evolution, again evident in \secref{subsec:literature-review-results-sos-characteristics}, indicates that many studies conceptualize SoTS as operating in stable settings despite the inherent instability of distributed and autonomous systems. The uneven adoption of standards described in \secref{subsec:literature-review-results-standards} shows progress toward interoperability but provides little assurance of reliable behaviour at runtime.

% Considered as a whole, the literature shows that although the structural and conceptual foundations of SoTS are becoming clearer, the operational capabilities needed to maintain reliable functioning are not yet developed. The primary goals that motivate SoTS, enhanced coordination, shared situational awareness, and informed decision-making, depend on system behaviour that remains correct and continuous even as conditions change. Yet current approaches offer limited support for ensuring that such behaviour is preserved when systems reconfigure, when the data they depend on is disrupted, or when operating environments become unstable. This gap represents a central limitation that must be addressed for SoTS to progress beyond prototype demonstrations and operate reliably in practice.


\subsection{Additional Reccomendations}
The reliability gap identified above is part of a broader set of limitations observed in the current SoTS landscape. The following additional recommendations synthesize further themes that emerged across the review and highlight structural, methodological, and standardization challenges that hinder the development of scalable SoTS.

\begin{table*}[htp]
\centering
\setlength{\tabcolsep}{1em}
\caption{Challenges}
\label{tab:challenges-table}
\footnotesize
\begin{tabular}{@{}p{4cm} l p{8cm}@{}}
\toprule
\textbf{Challenge} & \textbf{Frequency} & \textbf{Studies} \\
\midrule
\textbf{Operational Challenges} & \textbf{\maindatabar{60}} & \\
\;\;\corner{} Interoperability & \subdatabar{26} & \cite{acharya2023twins,alam2017c2ps,chen2018digital,dahmen2022modeling,dobie2024network,esterle2021digital,gollner2022collaborative,heithoff2023challenges,hofmeister2024cross-domain,jiang2022novel,jirsa2024use,kulkarni2019towards,larsen2024towards,li2022cognitive,lippi2023enabling,marah2023architecture,park2020digital,parri2019jarvis,pickering2023towards,pillai2023digital,samak2023autodrive,schluse2017experimentable,somma2023digital,vermesan2021internet,villalonga2021decision-making,vogel-heuser2021approach} \\
\;\;\corner{} Synchronization & \subdatabar{11} & \cite{acharya2023twins,altamiranda2019system,ashtaritalkhestani2019architecture,bertoni2022digital,coupaye2023graph-based,duan2023digital,esterle2021digital,li2022cognitive,monsalve2021novel,novak2022digitalized,pillai2023digital} \\
\;\;\corner{} Real-Time Constraints & \subdatabar{9} & \cite{becue2018cyberfactory,gill2022method,hofmeister2024cross-domain,hofmeister2024semantic,joseph2021aggregated,malayjerdi2022combined,park2020digital,priyanta2024is,zhang2021bi-level} \\
\;\;\corner{} Uncertainty & \subdatabar{8} & \cite{bellavista2023requirements,bertoni2022digital,clark2021chapter,coupaye2023graph-based,demir2023vertically-integrated,oquendo2019dealing,parri2021framework,wang2024construction} \\
\;\;\corner{} Emergent Behaviors & \subdatabar{7} & \cite{barden2022academic,chen2018digital,dahmen2022modeling,gil2024integrating,kruger2022towards,li2022cognitive,liu2020web-based} \\
\;\;\corner{} Cost & \subdatabar{6} & \cite{ehemann2023digital,gill2022method,hatakeyama2018systems,hatledal2020co-simulation,mavromatis2024umbrella,pickering2023towards} \\
\;\;\corner{} Data Interoperability & \subdatabar{6} & \cite{doubell2023digital,kruger2022towards,li2024comprehensive,mahoro2023articulating,park2020digital,somma2023digital} \\
\;\;\corner{} Lifecycle Management & \subdatabar{4} & \cite{altamiranda2019system,aziz2022empowering,esterle2021digital,heithoff2023challenges} \\
\;\;\corner{} Adoption & \subdatabar{4} & \cite{becue2018cyberfactory,demir2023vertically-integrated,gill2022method,pickering2023towards} \\
\;\;\corner{} Decision Making & \subdatabar{4} & \cite{alam2017c2ps,barden2022academic,clark2021chapter,zhang2021bi-level} \\
\;\;\corner{} Reconfiguration & \subdatabar{4} & \cite{clark2021chapter,kruger2022towards,oquendo2019dealing,redelinghuys2020six-layer} \\
\;\;\corner{} Processing Efficiency & \subdatabar{3} & \cite{ehemann2023digital,marah2023architecture,saraeian2022digital} \\
\textbf{Design Challenges} & \textbf{\maindatabar{33}} & \\
\;\;\corner{} Complexity & \subdatabar{12} & \cite{bao2024digital,dickopf2019holistic,duan2023digital,ehemann2023digital,gill2022method,lee2022simulation,malayjerdi2022combined,marah2023architecture,pillai2023digital,saraeian2022digital,schluse2017experimentable,zhang2022multi-scale} \\
\;\;\corner{} Lack Of Standards & \subdatabar{11} & \cite{acharya2023twins,binder2021utilizing,coupaye2023graph-based,dickopf2019holistic,gill2022method,hatledal2020co-simulation,hofmeister2024cross-domain,howard2021greenhouse,jirsa2024use,larsen2024towards,vogel-heuser2021approach} \\
\;\;\corner{} Compatibility With Legacy Systems & \subdatabar{7} & \cite{dobie2024network,ehemann2023digital,gill2022method,howard2021greenhouse,lippi2023enabling,liu2020web-based,lopez2023modeling} \\
\;\;\corner{} Regulatory Constraints & \subdatabar{3} & \cite{malayjerdi2022combined,mavromatis2024umbrella,wullink2024foundational} \\
\;\;\corner{} Lack Of Frameworks/Architectures & \subdatabar{3} & \cite{ashtaritalkhestani2019architecture,howard2021greenhouse,villalonga2021decision-making} \\
\;\;\corner{} Collaboration & \subdatabar{3} & \cite{barden2022academic,demir2023vertically-integrated,mavromatis2024umbrella} \\
\;\;\corner{} Sociotechnical Integration & \subdatabar{2} & \cite{folds2019digital,mavromatis2024umbrella} \\
\;\;\corner{} Knowledge Representation & \subdatabar{2} & \cite{gil2023modeling,vogel-heuser2021approach} \\
\textbf{Non-Functional Properties} & \textbf{\maindatabar{22}} & \\
\;\;\corner{} Scalability & \subdatabar{6} & \cite{bertoni2022digital,chavezbaliguat2023digital,clark2021chapter,howard2021greenhouse,pillai2023digital,vermesan2021internet} \\
\;\;\corner{} Reliability & \subdatabar{4} & \cite{altamiranda2019system,aziz2022empowering,hofmeister2024semantic,kutzke2021subsystem} \\
\;\;\corner{} Privacy & \subdatabar{4} & \cite{heininger2021capturing,heithoff2023challenges,stary2022privacy,vermesan2021internet} \\
\;\;\corner{} Usability & \subdatabar{3} & \cite{chavezbaliguat2023digital,mavromatis2024umbrella,wagner2023using} \\
\;\;\corner{} Fidelity & \subdatabar{3} & \cite{folds2019digital,potteiger2023live,saraeian2022digital} \\
\;\;\corner{} Safety & \subdatabar{2} & \cite{joseph2021aggregated,savur2019hrc-sos} \\
\;\;\corner{} Security & \subdatabar{2} & \cite{becue2018cyberfactory,dobie2024network} \\
\bottomrule
\end{tabular}
\end{table*}

%------------------------------------------------------------
\subsection*{Architectural Foundations Remain Underdeveloped}

The review shows that the architectural foundations of SoTS remain limited, particularly in their ability to support scalable and autonomous system compositions. As reported in \secref{subsec:literature-review-results-architecture}, acknowledged and directed SoTS dominate current work (\tabref{tab:sots-type-table}), reflecting a reliance on central coordination rather than distributed autonomy. Collaborative and virtual SoTS, which are more aligned with adaptive and multi-actor configurations, appear far less frequently. This imbalance corresponds with the challenges identified in \secref{subsec:literature-review-results-motivations} and \tabref{tab:challenges-table}, where increasing system complexity, unclear architectural boundaries, and persistent coordination difficulties were recurring concerns. Although some studies adopt modeling frameworks such as SysML or UML, these approaches rarely produce reusable architectures suited to cross-organizational or dynamically reconfigurable environments. Taken together, the results indicate that SoTS require stronger architectural foundations that explicitly support autonomy, coordination, and the integration of multiple twins at scale.

\begin{extrarecoframe}{Recommendation}
Develop explicit architectural specifications and reusable reference models to support scalable, autonomous, and dynamically coordinated SoTS.
\end{extrarecoframe}



%------------------------------------------------------------
\subsection*{Standard Adoption Is Growing but Fragmented}

The literature reveals increasing but uneven use of standards in SoTS, demonstrating progress toward interoperability but also exposing fragmentation across systems. As shown in \secref{subsec:literature-review-results-standards}, fewer than half of the studies reference recognized standards, and those that do tend to focus on data and communication layers, such as OPC~UA, the Asset Administration Shell, and RAMI~4.0 (\tabref{tab:standards-table}). Broader SoS-oriented frameworks appear inconsistently, and most studies use standards as general guidance rather than fully adopting them as integration constraints. This fragmented landscape limits cross-system compatibility, complicates multi-organizational deployment, and constrains the ability to compose heterogeneous twins into larger SoTS. These trends indicate a need for more comprehensive and aligned standardization efforts that address interfaces, semantics, coordination patterns, and lifecycle considerations.

\begin{extrarecoframe}{Recommendation}
Advance and align DT- and SoS-oriented standards to support interoperable, multi-system SoTS ecosystems.
\end{extrarecoframe}


%------------------------------------------------------------
\subsection*{Dynamic SoS Properties Are Weakly Represented}

Dynamic SoS characteristics remain insufficiently supported in current SoTS research. As observed in \secref{subsec:literature-review-results-sos-characteristics}, structural dimensions such as distribution and independence appear frequently, while dynamic dimensions such as reconfiguration and evolution are addressed far less often. Emergent behaviour, a defining aspect of SoS, is typically examined only in controlled or simulation-based settings, with limited attention to strong or adaptive forms of emergence. These results show that current SoTS implementations often rely on static coordination patterns, even though the distributed and autonomous nature of SoS makes runtime change unavoidable. Digital Twins are well positioned to address this limitation, as they offer sensing, prediction, and simulation capabilities that could support runtime detection and interpretation of changing conditions. However, the review indicates that these capabilities are underutilized in practice.

\begin{extrarecoframe}{Recommendation}
Use DT sensing and simulation capabilities to detect, interpret, and manage emergent and reconfigurable behaviour at runtime.
\end{extrarecoframe}


%------------------------------------------------------------
\subsection*{Research Maturity Remains Low}

The maturity of SoTS research remains low, limiting confidence in how current approaches perform under real conditions. As shown in \secref{subsec:literature-review-results-trl}, most studies fall within the lower Technology Readiness Levels (\tabref{tab:trl-table}) and rely on simulations or controlled laboratory experiments. Only one study reports operational deployment, and few examine SoTS behaviour under realistic or adversarial conditions. This limited empirical grounding constrains understanding of system performance, reliability, and adaptability, especially when scaled across organizational or physical boundaries. The review also highlights barriers related to regulatory requirements, integration with legacy systems, and multi-stakeholder coordination. Addressing these limitations requires stronger evidence obtained through field-based studies, case study evaluations, and validation in settings that reflect the complexity of real-world SoTS applications.

\begin{extrarecoframe}{Recommendation}
Increase empirical, field-based evaluations to establish stronger evidence for SoTS performance, reliability, and adaptability in real operational environments.
\end{extrarecoframe}

%------------------------------------------------------------



\section{Conclusion}
These takeaways show that SoTS research is advancing in scope but remains limited in standardization, dynamic capability, and empirical grounding. Addressing these gaps is essential for enabling SoTS to move from early-stage prototypes toward deployable systems capable of operating in uncertain and rapidly changing environments. The next chapter builds on these observations to develop an architectural approach that responds directly to the reliability challenges identified here while aligning with the broader needs for scalable, adaptable, and well-structured SoTS.
